\documentclass[11pt]{article}
\usepackage[margin=1.15in]{geometry}
\usepackage{amsmath,amssymb,amsthm,mathtools}
\usepackage{xcolor}
\usepackage[colorlinks=true,linkcolor=blue!60!black,citecolor=blue!60!black,urlcolor=blue!60!black]{hyperref}

\newtheorem{theorem}{Theorem}[section]
\newtheorem{proposition}[theorem]{Proposition}
\newtheorem{lemma}[theorem]{Lemma}
\newtheorem{corollary}[theorem]{Corollary}
\newtheorem{conjecture}[theorem]{Conjecture}
\theoremstyle{definition}
\newtheorem{definition}[theorem]{Definition}
\theoremstyle{remark}
\newtheorem{remark}[theorem]{Remark}

\newcommand{\C}{\mathbb{C}}
\newcommand{\R}{\mathbb{R}}
\newcommand{\Z}{\mathbb{Z}}
\newcommand{\w}{\omega}
\newcommand{\dg}{\dagger}
\newcommand{\ad}{\operatorname{ad}}
\newcommand{\ob}{\mathfrak o}
\newcommand{\kerad}{\Pi_{\ker\ad_{h_0}}}

\title{Interaction Obstructions and Resonant Breakdown\\
of the Positive Pais--Uhlenbeck Metric}
\author{GPT-5.6.sol\thanks{OpenAI model.  Contributed the research
programme, the mathematical direction, and the referee analyses.}
\and
Claude Fable 5\thanks{Anthropic model.  Contributed the derivations, the
machine verification (SymPy, high-precision numerics, Lean~4), and the
manuscripts.  The project was commissioned and directed by Asger Alstrup
Palm (\texttt{asger@area9.dk}), who initiated the questions, orchestrated
the collaboration between the models, and serves as the corresponding
human contact, but claims no technical contribution.}}
\date{}

\begin{document}
\maketitle

\begin{abstract}
The free unequal-frequency Pais--Uhlenbeck (PU) oscillator admits a
canonical positive pseudo-Hermitian metric, geometrically selected and
classified in a companion series.  This paper asks what happens to that
metric --- and to its Krein/ghost-parity alternative --- when an
interaction is switched on.  For the oscillator with the cubic vertex
inherited from a real $z^3$ interaction we construct the metric
deformation explicitly through third order and prove: the deformation
exists for generic frequencies; at the internal resonance
$\w_1=3\w_2$ it is obstructed at second order by an exact on-shell
conversion operator, and sufficiently excited resonant multiplets then
acquire complex-conjugate energy pairs at $O(\zeta^2)$ --- genuine
spectral PT-breaking, which excludes \emph{every} positive invariant
metric where it persists; at $\w_1=\tfrac32\w_2$ a third-order
obstruction appears; a transfer-lattice selection rule organizes the
resonance hierarchy; and the deformation coefficients diverge
geometrically, $R_n=O(\epsilon^{-3n/2})$, toward the Jordan boundary.
For the perfect-square field theory
$-\tfrac12\int(\Box\phi+\lambda(\partial\phi)^2)^2$, regulated by a mass
splitting, the cubic vertex is exactly the K\"all\'en polynomial
$\tfrac12\lambda_{\mathrm K}(p_1^2,p_2^2,p_3^2)$; an even-ghost
selection rule protects the positive metric at first order even above
the ghost-decay threshold, while the sectorwise ghost parity is
obstructed there by the real decay.  At second order the momentum
continuum turns the oscillator's isolated resonances into generic
scattering shells: the branch-changing process $H+L\to L+L$ obstructs
the pointed-sector positive metric, with the exact value
$401\sqrt6/39424$ at a rational kinematic point, hence on an open
subset of the shell.  An exact two-field rewriting,
$U=e^{\lambda\phi}$, $V=(\psi/\lambda)e^{-\lambda\phi}$, converts the
theory into $-\partial U\!\cdot\!\partial V+\tfrac{\lambda^2}2U^2V^2$
and exhibits the hidden ghost parity as the exact nonlinear involution
$U\leftrightarrow V$, which exchanges two oppositely oriented null
vacuum sectors and generates the $\Box^2$ Jordan dynamics about either
of them; the regulated branch parity has no bounded confluent limit in
one sector but converges on the doubled space to this exchange.  All
finite-dimensional and modewise identities are machine-verified.
\end{abstract}

\tableofcontents

\section{Introduction}\label{sec:intro}

A companion series \cite{paperI,paperII,paperIII,paperIV} analyzed the
\emph{free} fourth-order theories --- the PU oscillator, the fourth-order
scalar field, and linearized quadratic gravity --- and established a
three-level structure: one complex spectral covariance, several real
forms (involutions), and several completions.  The positive
pseudo-Hermitian completion of Bender and Mannheim
\cite{BM2008PRL,BM2008PRD} and the Krein/ghost-parity completion of
Bateman and Turok \cite{BT2026} are real forms of one free theory,
distinguished at Jordan degenerations.

This paper begins the interacting classification.  The organizing
object is the \emph{deformation problem} for the canonical positive
metric: given the free Dyson map $\rho_0$ with
$h_0=\rho_0H_0\rho_0^{-1}$ Hermitian, and a perturbation
$H_\zeta=H_0+\zeta V$, when does a deformed map
$\rho_\zeta=\exp[-\tfrac12(\zeta R_1+\zeta^2R_2+\cdots)]\rho_0$ keep
$h_\zeta=\rho_\zeta H_\zeta\rho_\zeta^{-1}$ Hermitian?  Order by order
this is a cohomological question: the order-$n$ equation is
$[h_0,R_n]=S_n$ with a source built from lower orders, solvable iff the
kernel projection $\kerad S_n$ --- the \emph{obstruction class} ---
vanishes.  A parallel problem governs the Krein side: deform the ghost
parity $\kappa_0$ so that $[\kappa_\zeta,H_\zeta]=0$,
$\kappa_\zeta^2=1$.

\paragraph{Results (oscillator).}
For the unequal-frequency PU oscillator with the cubic vertex
$V=-iy^3$ inherited from a real $z^3$ interaction:
\begin{enumerate}
\item the first-order deformation exists for all $\w_1>\w_2>0$, with an
explicit Weyl-ordered cubic $R_1$ (Section~\ref{sec:first});
\item the second-order deformation exists for generic (nonresonant)
frequencies, with the assembled $h_\zeta$ Hermitian through
$O(\zeta^2)$ (Section~\ref{sec:second});
\item at the internal $3{:}1$ resonance $\w_1=3\w_2$ the second-order
class is nonzero and exactly
$\ob_+^{(2)}=\tfrac{27\sqrt3}{320\w_2^4}\,(a_1a_2^{\dg3}-a_1^{\dg}a_2^3)$,
the on-shell $1\leftrightarrow3$ conversion, unremovable by any
first-order freedom (Section~\ref{sec:threeone});
\item the obstruction is spectral, not merely metric: the
$E=27\w_2$ resonant multiplet develops a complex-conjugate energy pair
at $O(\zeta^2)$, verified both in degenerate perturbation theory and by
exact diagonalization of the truncated Hamiltonian; the broken-PT
region has width $|\w_1-3\w_2|=O(\zeta^2)$
(Section~\ref{sec:spectral});
\item a third-order obstruction appears at $\w_1=\tfrac32\w_2$; a
transfer-lattice and degree-parity selection rule organizes the
candidate loci, and the verified pattern is that the coprime ratio
$p{:}q$ first obstructs at order $p+q-2$ when the mode-2 transfer $p$
is odd (Section~\ref{sec:selection});
\item toward the Jordan boundary $\epsilon=\tfrac12(\w_1-\w_2)\to0$ the
deformation diverges geometrically, $R_n=O(\epsilon^{-3n/2})$
(Section~\ref{sec:jordanscaling}).
\end{enumerate}

\paragraph{Results (perfect-square field theory).}
For the interacting theory
$S=-\tfrac12\int(\Box\phi+\lambda(\partial\phi)^2)^2$, regulated by a
mass splitting $\delta=m_1^2-m_2^2>0$:
\begin{enumerate}
\setcounter{enumi}{6}
\item the symmetrized cubic vertex is exactly
$\tfrac12\lambda_{\mathrm K}(p_1^2,p_2^2,p_3^2)$, the K\"all\'en
polynomial; it vanishes on three massless shells and on one generalized
Jordan leg, and couples Jordan modes in pairs
(Section~\ref{sec:kallen});
\item at first order an \emph{even-ghost selection rule} makes the
positive obstruction vanish identically --- below and above the
ghost-decay threshold $m_1=2m_2$ --- while the sectorwise ghost parity
$(-1)^{N_{\mathrm{ghost}}}$ is obstructed above threshold by the real
decay, with class proportional to
$\lambda_{\mathrm K}(m_1^2,m_2^2,m_2^2)=m_1^2(m_1^2-4m_2^2)$
(Section~\ref{sec:firstfield});
\item at second order the momentum continuum makes the resonances
generic: the branch-changing scattering $H+L\to L+L$ is exactly energy
conserving at open sets of momenta for \emph{every} split mass pair,
and the positive obstruction is nonzero there --- with the exact value
$401\sqrt6/39424$ at the rational center-of-mass point
$m_H=\tfrac32m_L$ (Section~\ref{sec:sector});
\item the theory admits an exact two-field rewriting
$\mathcal L=-\partial U\!\cdot\!\partial V+\tfrac{\lambda^2}2U^2V^2$
with $U=e^{\lambda\phi}$, $V=(\psi/\lambda)e^{-\lambda\phi}$; the
hidden ghost parity is the exact nonlinear involution
$U\leftrightarrow V$, which exchanges the two null vacuum branches and
generates the $\Box^2$ Jordan dynamics about either of them
(Section~\ref{sec:twofield});
\item the regulated branch parity has no bounded confluent limit inside
one Jordan sector, but converges, on the doubling by the oppositely
oriented sector, to the exact sector-exchange involution
(Section~\ref{sec:confluent}).
\end{enumerate}

\paragraph{The central statement.}
\begin{quote}
\emph{Positive pseudo-Hermitian metrics can be deformed locally away
from resonances, but interactions generate cohomological obstructions
whenever branch-changing processes become exactly energy conserving.
In finite systems these occur on isolated rational resonance loci and
can produce complex spectra.  In field theory the momentum continuum
makes the corresponding scattering shells generic.  The exact massless
perfect-square theory nevertheless possesses a nonperturbative doubled
involution --- not as a parity within either pointed Jordan sector, but
as an exchange between two oppositely oriented sectors.}
\end{quote}

\paragraph{Bookkeeping.}
Two expansion parameters must not be conflated.  The \emph{deformation
order} is counted by a formal parameter $\zeta$:
$H(\zeta)=H_0+\zeta H_3+\zeta^2H_4$, defining $R_1,R_2,\dots$ and the
obstruction orders.  The perfect-square coupling $\lambda$ is a fixed
parameter: in the two-field frame the Jordan quadratic structure and
the vertex coefficients carry $g=\lambda^2$, while the nonlinear field
map itself depends on $\lambda$.  In the oscillator sections
$\zeta$ multiplies the cubic vertex and coincides with the coupling.

\paragraph{Verification.}
Every finite-dimensional, modewise, and distributional identity in this
paper is machine-verified: the first-order theory and second-order
recursion (checks ID1--ID10), selection rules and third order
(SR1--SR3, O3a--O3e), spectral PT-breaking and the K\"all\'en vertex
(PT1--PT4, PS1--PS2 of \texttt{verify\_pt\_breaking.py}), the
first-order field theory (PS-A--PS-H), the two-field rewriting and
exact involution (TF1--TF7), the sectorwise second-order obstruction
(SO1--SO7), and the exact kinematic point and confluent parity
(HX1--HX3), in the repository accompanying the series (SymPy 1.14.0,
NumPy 2.4.4, Python 3.14; release tag \texttt{paper5-v1.0} fixes the
commit).  Appendix~\ref{app:checks} maps checks to theorems.
Operator-domain statements (unbounded cubic interactions, formal
perturbation series) are flagged where they occur.

\paragraph{Conventions.}
Oscillator conventions follow \cite{paperI,paperII}: normalized
coordinates $V=(x,y,p,q)$, mode 1 $=(x,p)$ at frequency $\w_1$, mode 2
$=(y,q)$ at $\w_2$, positive normal form
$h_0=\tfrac12[\w_1\w_2p^2+(\w_1/\w_2)x^2+(\w_2/\w_1)q^2+\w_1\w_2y^2]$,
rapidity $r=\log\frac{\w_1+\w_2}{\w_1-\w_2}$,
$c=\cosh\tfrac r2=\w_1/\sigma$, $s=\sinh\tfrac r2=\w_2/\sigma$,
$\sigma=\sqrt{\w_1^2-\w_2^2}$.  All operator computations use exact
Weyl-symbol Moyal calculus (finite for polynomials; commutators with
the quadratic $h_0$ reduce to $i\times$ Poisson bracket; cubic--cubic
commutators include the $\Lambda^3$ term).

\section{The deformation complex}\label{sec:complex}

\begin{definition}[Deformation problem and obstruction classes]
Let $\rho_0$ be the canonical free Dyson map,
$h_0=\rho_0H_0\rho_0^{-1}=h_0^{\dg}$, and let
$v=\rho_0V\rho_0^{-1}$ be the transported vertex.  Seek
$\rho_\zeta=\exp[-\tfrac12\sum_n\zeta^nR_n]\rho_0$ with
$R_n=R_n^{\dg}$ such that
$h_\zeta=\rho_\zeta H_\zeta\rho_\zeta^{-1}$ is Hermitian order by
order.  The first two equations are
\begin{equation}\label{eq:defeq}
[h_0,R_1]=v^{\dg}-v,
\qquad
[h_0,R_2]=(v_2^{\dg}-v_2)+\tfrac12[R_1,\,v+v^{\dg}],
\end{equation}
where $v_2$ is the transported $O(\zeta^2)$ vertex when present.  With
$h_0=\sum_EE\,P_E$, the equation $[h_0,R]=S$ has a Hermitian solution
iff $P_ESP_E=0$ for every $E$; the obstruction class is
$\ob_+ = \kerad S$, and the solution is
$R=\sum_{E\neq E'}P_ESP_{E'}/(E-E')$ plus an arbitrary Hermitian
element of the commutant.  The parallel Krein problem deforms the ghost
parity: $[H_0,K_1]=[\kappa_0,V]$, $\{\kappa_0,K_1\}=0$, with
$\ob_K=\Pi_{\ker\ad_{H_0}}[\kappa_0,V]$.
\end{definition}

Anti-Hermitian (unitary) components of the generator are free and do
not affect the Hermiticity conditions; additions to $R_n$ from the
commutant of $h_0$ are the usual metric ambiguity.  ``Gauge
independence'' of an obstruction class means invariance under both.

\section{The cubic PU oscillator at first order}\label{sec:first}

The real interaction $z^3$ continues, under the PU rotation $z=iy$, to
$V=-iy^3$.  The transported vertex is $v=-i(cy+isp)^3$ (the free Dyson
map acts by $\rho_0y\rho_0^{-1}=cy+isp$), whence
\begin{equation}
v^{\dg}-v=2ic\,(c^2y^3-3s^2yp^2).
\end{equation}

\begin{theorem}[First-order deformation]\label{thm:first}
For every $\w_1>\w_2>0$ the first-order equation is solved by the
Hermitian Weyl-ordered cubic
\begin{equation}
R_1=\mathcal W\!\left[
\frac{2c^3}{\w_1\w_2}qy^2
+\frac{4c^3}{3\w_1^3\w_2}q^3
-\frac{6cs^2(2\w_1^2-\w_2^2)}{\w_1\w_2(4\w_1^2-\w_2^2)}p^2q
+\frac{12cs^2\w_1}{\w_2(4\w_1^2-\w_2^2)}pxy
-\frac{12cs^2\w_1}{\w_2^3(4\w_1^2-\w_2^2)}qx^2
\right],
\end{equation}
with no singular denominator in the physical region
($4\w_1^2-\w_2^2>0$); the equivalent Hermitian interaction is
$h_1=\tfrac12(v+v^{\dg})=s(3c^2y^2p-s^2p^3)$.  Near the Jordan
boundary, $c\sim s\sim\sqrt\w/(2\sqrt\epsilon)$ and
$R_1=O(\epsilon^{-3/2})$, against the logarithmic divergence of the
free generator $Q_0$.
\end{theorem}

\section{Second order: generic existence}\label{sec:second}

\begin{theorem}[Generic second-order deformation]\label{thm:secondgen}
For incommensurate frequencies the second-order obstruction
$\ob_+^{(2)}=\kerad\tfrac12[R_1,v+v^{\dg}]$ vanishes; the solved $R_2$
is Hermitian, and the assembled
$h_\zeta=e^{-X/2}\star(h_0+\zeta v)\star e^{X/2}$,
$X=\zeta R_1+\zeta^2R_2$, is Hermitian through $O(\zeta^2)$ ---
verified end to end in exact Moyal calculus, including the
$\tfrac18[[h_0,R_1],R_1]$ term and the $\Lambda^3$ contribution that
pure Poisson-bracket reasoning would miss.  $R_2$ carries the
resonance denominators $(\w_1-\w_2)$ and $(\w_1-3\w_2)$; no others can
vanish for $\w_1>\w_2>0$.
\end{theorem}

\section{The exact $3{:}1$ obstruction}\label{sec:threeone}

\begin{theorem}[Second-order obstruction at $3{:}1$]\label{thm:threeone}
At $\w_1=3\w_2$ the first-order deformation still exists, but the
second-order class is nonzero and exact:
\begin{equation}
\ob_+^{(2)}
=\frac{27\sqrt3}{320\,\w_2^4}
\left(a_1a_2^{\dg3}-a_1^{\dg}a_2^3\right),
\end{equation}
the anti-Hermitian on-shell conversion of one mode-1 quantum into three
mode-2 quanta ($\w_1=3\w_2$ makes the process exactly energy
conserving).  The class is gauge independent: it is unchanged under
$R_1\to R_1+\{N_1,N_2,N_1N_2\}$-type commutant additions, under
additions of the on-shell quintic kernel operators to $R_2$, and under
anti-Hermitian first-order generators ---
$\kerad[G,v+v^{\dg}]=\kerad[v-v^{\dg},G]=0$ for all such $G$.
\end{theorem}

There is therefore no analytic second-order deformation of the
canonical positive metric at $3{:}1$, within the Weyl-algebra class,
continuously connected to the free metric.  The physical meaning: the
interaction contains a process that is exactly on shell, and changing
the metric cannot rotate it away.

\section{Spectral PT-breaking}\label{sec:spectral}

The obstruction is not an artifact of the deformation scheme: it enters
the spectrum.

\begin{theorem}[Complex pair in the $E=27\w_2$ multiplet]
\label{thm:complexpair}
At $\w_1=3\w_2$ (set $\w_2=1$), the second-order degenerate
perturbation matrix on the ten-dimensional shell
$\{|j,27-3j\rangle\}_{j=0}^9$ is real tridiagonal with antisymmetric
off-diagonal $K_{j,j+1}=-\tfrac{27\sqrt3}{640}\sqrt{(j{+}1)n(n{-}1)(n{-}2)}$,
$n=27-3j$, and explicit diagonal shifts.  Its spectrum consists of
eight real eigenvalues and one complex-conjugate pair
\begin{equation}
\kappa_\pm=2.44880738219938\pm0.224586593808125\,i .
\end{equation}
The pair is genuinely spectral: exact diagonalization of the truncated
$h_0+\zeta v$ at $\zeta=0.02$ exhibits it with
$\operatorname{Im}E=\pm0.2246\,\zeta^2$ to $0.5\%$, stable across
cutoffs.  The lowest resonant doublet $\{|1,0\rangle,|0,3\rangle\}$
remains real: a nonzero metric obstruction does not automatically break
the lowest pair.
\end{theorem}

\begin{corollary}[Strength of the conclusion]
Wherever the complex pair persists in the interacting spectrum, no
positive-definite invariant metric exists --- analytic in $\zeta$ or
not --- since a positive metric would make $h_\zeta$ similar to a
self-adjoint operator.  (The statement is perturbative-spectral: the
cubic interaction is unbounded and an operator-theoretic
analytic-family argument would be needed for the untruncated theory.)
\end{corollary}

\begin{proposition}[Broken-PT tongues]
Detuning $\delta_r=\w_1-3\w_2$ adds $j\,\delta_r$ to the shell
diagonal; the spectrum is complex for
$|\delta_r|<\delta_c=38.4151\ldots\times\zeta^2$ (exactly linear in
$\zeta^2$; the constant is multiplet-specific).  An order-$n$ resonance
correspondingly produces an $O(\zeta^n)$ tongue.
\end{proposition}

\section{Third order and the selection rules}\label{sec:selection}

\begin{theorem}[Transfer-lattice selection rule]\label{thm:selection}
Grade monomials by energy transfer
$(\Delta n_1,\Delta n_2)$, so $[h_0,M]=[(\Delta n_1)\w_1
+(\Delta n_2)\w_2]\,M$.  Order-$n$ objects have transfers in the
$n$-fold sumset of the vertex transfer set, total degree $\le n+2$
(each commutator reduces degree by at least two), and degree parity
$\equiv n\bmod2$.  The order-$n$ obstruction is therefore supported on
the finite candidate sets: $\{2{:}1\}$ at order 1, $\{3{:}1\}$ at
order 2, $\{3{:}2,\,2{:}1,\,4{:}1\}$ at order 3, with $5{:}1$ first
reachable at order 4.  The lattice gives candidates; coefficients
decide.
\end{theorem}

\begin{theorem}[Third order]\label{thm:third}
At generic frequencies the third-order obstruction vanishes and $R_3$
is Hermitian with odd total degrees $\le5$ (confirming the parity
rule).  At $\w_1=\tfrac32\w_2$ the class is nonzero and exact,
\begin{equation}
\ob_+^{(3)}
=-\frac{117\sqrt{30}}{1120}\,i\,
\big(a_1^2a_2^{\dg3}+a_1^{\dg2}a_2^3\big),
\end{equation}
the on-shell $2\leftrightarrow3$ conversion ($2\w_1=3\w_2$), gauge
independent in the sense of Theorem~\ref{thm:threeone}.  The candidates
$2{:}1$ (orders 1 and 3) and $4{:}1$ (order 3) have vanishing
coefficients.
\end{theorem}

\begin{conjecture}[Resonance hierarchy]\label{conj:hierarchy}
The coprime ratio $\w_1/\w_2=p{:}q$ first obstructs at order $p+q-2$
when the mode-2 transfer $p$ is odd.  (Verified: $3{:}1$ at order 2,
$3{:}2$ at order 3.  Predicted: $5{:}1$ at order 4, $5{:}3$ at
order 6.  Every obstruction found so far carries odd mode-2 transfer;
the mechanism excluding even transfers is open.)  The resonance set is
an expanding hierarchy, not (so far) a dense set.
\end{conjecture}

\section{The Jordan divergence hierarchy}\label{sec:jordanscaling}

\begin{theorem}[Geometric divergence]\label{thm:jordanscaling}
Near the Jordan boundary $\w_{1,2}=\w\pm\epsilon$,
\begin{equation}
Q_0=O(\log\epsilon^{-1}),\qquad
R_1=O(\epsilon^{-3/2}),\qquad
R_2=O(\epsilon^{-3}),\qquad
R_3=O(\epsilon^{-9/2}),
\end{equation}
the last three measured as clean decade exponents ($-1.50$, $-3.00$,
$-4.49$); the second order shows \emph{no} small-denominator
enhancement from $\w_1-\w_2$ (those source components are suppressed).
\end{theorem}

\begin{conjecture}
$R_n=O(\epsilon^{-3n/2})$, so the deformation's radius of convergence
in $\zeta$ collapses as $\zeta_c\sim\epsilon^{3/2}$ toward the Jordan
boundary.  (Fixed-order scaling is established; a norm bound on the
coefficients would upgrade this to a theorem.)
\end{conjecture}

\section{The perfect-square vertex}\label{sec:kallen}

The interacting theory is
$S=-\tfrac12\int(\Box\phi+\lambda(\partial\phi)^2)^2$
\cite{BT2026}, with cubic vertex $V_3\propto\Box\phi(\partial\phi)^2$
and quartic $V_4\propto\tfrac12(\partial\phi)^4$.

\begin{proposition}[K\"all\'en form]\label{prop:kallen}
On $p_1+p_2+p_3=0$ the symmetrized cubic vertex is exactly
\begin{equation}
\mathcal V_3
=p_1^2(p_2\!\cdot\!p_3)+p_2^2(p_1\!\cdot\!p_3)+p_3^2(p_1\!\cdot\!p_2)
=\tfrac12\lambda_{\mathrm K}(p_1^2,p_2^2,p_3^2),
\end{equation}
with $\lambda_{\mathrm K}(x,y,z)=x^2+y^2+z^2-2xy-2xz-2yz$.
Consequently: $\mathcal V_3$ vanishes on three ordinary massless
shells; it vanishes with one generalized Jordan leg
($\nabla\lambda_{\mathrm K}|_0=0$); it can be nonzero with two
(Hessian entries $2$, $-2$): \emph{the cubic vertex couples the Jordan
sector in pairs}.  On the regulated decay channel,
$\lambda_{\mathrm K}(m_1^2,m_2^2,m_2^2)=m_1^2(m_1-2m_2)(m_1+2m_2)$:
the vertex vanishes exactly at the threshold $m_1=2m_2$ and turns on
continuously above it.
\end{proposition}

\section{First order in the regulated field theory}\label{sec:firstfield}

Regulate by $P_\delta=(\Box+m_1^2)(\Box+m_2^2)$,
$\delta=m_1^2-m_2^2>0$.  In the mode expansion every leg is on its
branch shell, so on the kernel of $\ad_{h_0}$ the full four-momentum is
conserved and the vertex coefficient is
$\tfrac12\lambda_{\mathrm K}$ of the branch masses.  The transported
field mode, derived from the free Dyson transport and verified against
the spectral Wightman function of \cite{paperIII}, is
\begin{equation}
\phi'(k)=\frac{i}{\sqrt{2\delta}}
\left[\frac{a_2+a_2^{\dg}}{\sqrt{\w_2}}
+\frac{a_1-a_1^{\dg}}{\sqrt{\w_1}}\right],
\qquad
\sigma=\sqrt{\w_1^2-\w_2^2}=\sqrt\delta\ \text{($k$-independent)}.
\end{equation}

\begin{theorem}[First-order protection and Krein decay obstruction]
\label{thm:firstfield}
(i) \emph{Even-ghost selection rule}: the transported reality factors
make $v^{\dg}-v$ vanish on every odd-ghost-count monomial.  Even-ghost
shells ($2\to22$, $11\to2$, \dots) are kinematically closed for all
$m_1>m_2>0$.  Hence $\ob_+^{(1)}=0$ identically --- below \emph{and}
above the decay threshold.
(ii) The sectorwise ghost parity $\kappa_0=(-1)^{N_{\mathrm{ghost}}}$
has $\ob_K^{(1)}=0$ below threshold, and above it a nonzero class
supported on the open $1\to22$ shell, equal there to
$\lambda_{\mathrm K}(m_1^2,m_2^2,m_2^2)$ times the leg normalizations:
a really decaying ghost has no conserved sectorwise parity, and the
protection is exactly coextensive with $\lambda_{\mathrm K}=0$
(threshold and massless point).
(iii) Scaling: at massive confluence both $R_1$ and $K_1$ diverge as
$\delta^{-3/2}$; on massless Jordan paths $m_2^2=\alpha\delta$ the
positive $R_1$ scales as $\delta^{-1/2}$ with path-independent power
(the selection rule removes the dangerous collinear term), while the
Krein $K_1$ is $\delta^{-1/2}$ generically and $\delta^{-3/2}$ on the
exceptional path $\alpha=2/7$.  (These are split-basis statements; the
basis itself degenerates at $\delta=0$, and a basis-independent
formulation on confluent states is deferred.)
\end{theorem}

\section{The exact two-field rewriting}\label{sec:twofield}

\begin{theorem}[Two-field form and exact involution]\label{thm:twofield}
Introducing the auxiliary field $\psi$ and the exact nonlinear change
of variables $U=e^{\lambda\phi}$,
$V=(\psi/\lambda)e^{-\lambda\phi}$ (unit pointwise Jacobian),
\begin{equation}
\mathcal L
=-\partial_\mu U\,\partial^\mu V+\frac{\lambda^2}{2}U^2V^2
=-\tfrac12(\partial X)^2+\tfrac12(\partial Y)^2
+\frac{\lambda^2}{8}(X^2-Y^2)^2,
\end{equation}
with $X,Y=(U\pm V)/\sqrt2$: the interaction is built from the
$O(1,1)$ invariant.  The exchange $U\leftrightarrow V$ is an exact
symmetry; in the original variables it is the nonlinear involution
\begin{equation}
\kappa_0:\ (\phi,\psi)\ \longmapsto\
\Big(\tfrac1\lambda\log\tfrac{\psi}{\lambda}-\phi,\ \psi\Big),
\end{equation}
invisible at any finite perturbative order in $(\phi,\psi)$ ---
geometrically, the transition map between the two charts
$\phi_A=\tfrac1\lambda\log U$ and $\phi_B=\tfrac1\lambda\log V$ on the
extended configuration space.
\end{theorem}

\begin{proposition}[Null sectors and interaction-generated Jordan
dynamics]\label{prop:sectors}
The vacuum manifold is $UV=0$: two null branches.  The perturbative
vacuum $(U,V)=(1,0)$ is a nonzero null background;
$\kappa_0$ maps it to $(0,1)$ --- it exchanges pointed sectors and is
not a particle parity within either.  Expanding about $(1,0)$:
\begin{equation}
\mathcal L^{(2)}=-\partial u\,\partial v+\tfrac{g}2v^2,\quad
\mathcal L^{(3)}=g\,uv^2,\quad
\mathcal L^{(4)}=\tfrac{g}2u^2v^2,\qquad g=\lambda^2,
\end{equation}
with equations $\Box v=0$, $\Box u=gv$: the exact $\Box^2$ Jordan pair,
the off-diagonal supplied by the interaction coupling.  About $(0,1)$
the mirror expansion has $u\leftrightarrow v$: the oppositely oriented
chain.  A chain-preserving involution commuting with a single Jordan
block is $\pm I$ (companion lemma), so the exchange structure is the
only possible realization of ghost parity --- as the exact
$\kappa_0$ indeed exhibits.  The interactions in this frame are pure
potentials: $H_{\mathrm{int}}=-L_{\mathrm{int}}$ exactly, with no
Ostrogradsky--Legendre corrections (in the $(\phi,\psi)$ frame the
quartic contains time derivatives and the canonical
second-order vertex differs from $-L_4$).
\end{proposition}

\section{The sectorwise second-order obstruction}\label{sec:sector}

Regulate the sector-$A$ theory by
$\mathcal L_0=-\partial u\partial v-\mu^2uv+\tfrac g2v^2
+\tfrac\epsilon2u^2$ (masses $m_\pm^2=\mu^2\pm\sqrt{\epsilon g}$; the
regulator necessarily breaks $\kappa_0$), quantize in the positive
frame via the branch eigenvectors ($v=\rho_bu$ on branch $b$,
$\rho_\pm=\mp\delta/2$), and use the exact matrix-element form of the
second-order source:
\begin{equation}\label{eq:sourceME}
\langle\mathrm{out}|\,\mathcal S^{(2)}_+\,|\mathrm{in}\rangle
=\langle v_2^{\dg}-v_2\rangle
+\sum_n\frac{\langle\mathrm{out}|v_1^{\dg}|n\rangle
\langle n|v_1^{\dg}|\mathrm{in}\rangle
-\langle\mathrm{out}|v_1|n\rangle\langle n|v_1|\mathrm{in}\rangle}
{E-E_n},
\end{equation}
which follows from
$(v^{\dg}-v)(v+v^{\dg})+(v+v^{\dg})(v^{\dg}-v)=2(v^{\dg}v^{\dg}-vv)$.

\begin{theorem}[Generic sectorwise obstruction]\label{thm:sector}
For the branch-changing scattering $H+L\to L+L$ --- kinematically open
at suitable momenta for \emph{every} $m_1>m_2$ --- the source element
\eqref{eq:sourceME} is nonzero.  At the rational center-of-mass point
$m_H=\tfrac32m_L$ (both incoming particles at rest, outgoing momenta
$\pm\tfrac34m_L$, all shell conditions exactly satisfied) its exact
value is
\begin{equation}
\boxed{\ \mathcal M_{\mathrm{obs}}=\frac{401\sqrt6}{39424}\ }
\end{equation}
(with $m_L=4$ in lattice units), the quartic contact piece vanishing
exactly.  Since the source is analytic on the nonsingular part of the
shell, the obstruction is nonzero on an open subset: \emph{the
second-order positive deformation of the pointed sector is generically
obstructed}.  The computation is truncation-complete (tree-level
exchange: the $s,t,u$ channels close on the momentum set; verified by
cutoff independence) and reproduced at three independent tuned
kinematic points.
\end{theorem}

\begin{remark}[Oscillator versus field theory]
\begin{center}
\begin{tabular}{c|c}
PU oscillator & field theory\\
\hline
discrete energy spectrum & continuous momentum spectrum\\
isolated rational resonance & generic scattering shell\\
$3{:}1$ first appears at order 2 & $H+L\to L+L$ appears at order 2\\
finite resonant multiplet & continuum of degenerate in/out states
\end{tabular}
\end{center}
The continuum turns the oscillator's discrete resonance obstruction
into a generic scattering-shell obstruction.  The exact doubled
$\kappa_0$ does not cure it: it maps the obstructed sector-$A$
construction to the (equally obstructed) mirror sector-$B$
construction.
\end{remark}

\begin{remark}[Strength of the field-theory statement]
Theorem~\ref{thm:sector} establishes a perturbative cohomological
obstruction: there is no analytic second-order deformation of the
canonical positive metric, in the specified regulated Fock
representation and operator class, for a pointed sector.  Unlike the
oscillator case (Theorem~\ref{thm:complexpair}), we have not yet
exhibited complex spectra here; the analogous upgrades would be complex
finite-volume levels, non-real resonance poles, or failure of
pseudo-unitarity of the pointed-sector $S$-matrix.  The distinction ---
analytic metric obstruction versus absolute nonexistence of all
positive metrics --- is maintained throughout.
\end{remark}

\section{The confluent parity theorem}\label{sec:confluent}

\begin{theorem}[Confluent limit of the regulated parity]
\label{thm:confluent}
Let $P_\delta b_\pm=\pm b_\pm$ be the regulated branch ghost parity and
$c=(b_++b_-)/2$, $d=(b_+-b_-)/\delta$ the confluent Jordan basis.  Then
\begin{equation}
P_\delta=
\begin{pmatrix}0&\delta/2\\ 2/\delta&0\end{pmatrix}
\quad\text{in the }(c,d)\text{ basis:}
\end{equation}
$P_\delta^2=1$ but $P_\delta$ has no bounded limit as $\delta\to0$ ---
the regulated ghost parity cannot converge as a chain-preserving parity
inside one Jordan sector.  On the doubled space with the
\emph{oppositely oriented} confluent identification in the mirror
sector ($c_B=(b_+-b_-)/2$, $d_B=(b_++b_-)/\delta$), the cross parity
$b_\pm^A\leftrightarrow\pm b_\pm^B$ equals, exactly and
$\delta$-independently, the sector exchange
$(c_A,d_A)\leftrightarrow(c_B,d_B)$.  The regulated branch parity thus
converges, after doubling by the oppositely oriented Jordan sector, to
the exact involution $U\leftrightarrow V$.
\end{theorem}

This connects the split Krein algebra directly to the exact nonlinear
symmetry, and explains the earlier observations from one statement: the
sectorwise $(-1)^{N_{\mathrm{ghost}}}$ is broken by real decay and has
no confluent limit, while the exact parity survives --- on the doubled
state space.

\section{Discussion}\label{sec:discussion}

\paragraph{The phase picture.}
Away from resonances the positive pseudo-Hermitian metric deforms
locally, at first and second order, with explicit generators.  The
failures are structural, not accidental: (i) on-shell branch-changing
processes produce gauge-independent cohomological obstructions ---
isolated rational loci for the oscillator, generic scattering shells in
the field theory; (ii) at the oscillator resonances the obstruction is
spectral, producing complex-conjugate pairs in excited multiplets and
$O(\zeta^2)$ broken-PT tongues; (iii) independent of resonances, the
deformation loses uniformity toward Jordan degeneration,
$R_n=O(\epsilon^{-3n/2})$.  The Krein side is not simply the survivor:
its sectorwise parity is broken by real ghost decay above threshold.
What survives is subtler --- an exact, nonperturbative, sector-exchanging
involution of the massless perfect-square theory, which the confluent
parity theorem identifies as the $\delta\to0$ limit of the regulated
branch parity on the doubled state space.

\paragraph{Symmetry enhancement at the boundary.}
The massless perfect-square point is not the continuous limit of the
split theory with parameters set to zero: the K\"all\'en vertex
degenerates there, the dangerous cubic amplitudes vanish, the doubled
$O(1,1)$ structure emerges, and the Jordan dynamics is generated by the
interaction about its own null vacua.  This mirrors the conformal
boundary of the gravity companion \cite{paperIV}, where a Weyl gauge
symmetry emerges at a singular locus and the boundary theory must be
quotiented and completed differently.

\paragraph{Superselection (open).}
Whether the two pointed sectors admit cat-like parity vacua
$|0_\pm\rangle=(|0_A\rangle\pm|0_B\rangle)/\sqrt2$, or are superselected
in infinite volume (spontaneously broken $\kappa_0$ in each pointed
completion), requires finite-volume overlap and zero-mode analysis:
the vacuum manifold is null, the kinetic form indefinite, and the
original chart contains only one branch.  The present results are
stated so as not to depend on the resolution.

\paragraph{Predictions and next steps.}
The hierarchy conjecture predicts a $5{:}1$ obstruction at fourth order
(an $O(\zeta^4)$ tongue); only the kernel projection at
$\w_1=5\w_2$ is needed to test it.  For the field theory, the
basis-independent confluent-state evaluation of $R_1$ will decide
whether the $\delta^{-1/2}$ split-basis divergence is physical.  The
gravitational extension --- the cubic Einstein--Weyl vertices $hhM$,
$hMM$, $MMM$ against the covariant real forms of \cite{paperIV} --- is
the natural target after that; the present results sharpen its
expected shape: pointed positive completions obstructed by
ghost-branch-changing graviton processes, with any surviving structure
living on a doubled completion.

\begin{thebibliography}{12}

\bibitem{BM2008PRL}
C.~M.~Bender and P.~D.~Mannheim,
\emph{No-ghost theorem for the fourth-order derivative Pais--Uhlenbeck
oscillator model},
Phys.\ Rev.\ Lett.\ \textbf{100} (2008) 110402.

\bibitem{BM2008PRD}
C.~M.~Bender and P.~D.~Mannheim,
\emph{Exactly solvable PT-symmetric Hamiltonian having no Hermitian
counterpart},
Phys.\ Rev.\ D \textbf{78} (2008) 025022, arXiv:0804.4190.

\bibitem{BT2026}
S.~Bateman and N.~Turok,
\emph{Escape from Ostrogradsky via hidden ghost parity},
arXiv:2607.00096.

\bibitem{Mostafazadeh2005}
A.~Mostafazadeh,
\emph{Metric operator in pseudo-Hermitian quantum mechanics and the
imaginary cubic potential},
J.\ Phys.\ A \textbf{39} (2006) 10171, arXiv:quant-ph/0508195.

\bibitem{paperI}
\emph{Canonical positive symplectic diagonalization of the
Pais--Uhlenbeck oscillator}, companion paper (Paper~I) and verification
repository, 2026.

\bibitem{paperII}
\emph{The Pais--Uhlenbeck metric as a minimum-distortion principle, and
the representation problem for the fourth-order field}, companion paper
(Paper~II), 2026.

\bibitem{paperIII}
\emph{The universal vacuum of the fourth-order scalar field: metric
orbits, Fock sectors, and the Krein boundary}, companion paper
(Paper~III), 2026.

\bibitem{paperIV}
\emph{Gauge reduction and the completion problem in fourth-order
gravity: PU pairing, covariant real forms, and the conformal Jordan
boundary}, companion paper (Paper~IV), 2026.

\bibitem{Kato}
T.~Kato,
\emph{Perturbation Theory for Linear Operators},
Springer, 1966.

\bibitem{Heiss2012}
W.~D.~Heiss,
\emph{The physics of exceptional points},
J.\ Phys.\ A \textbf{45} (2012) 444016.

\end{thebibliography}

\appendix

\section{Verification checks and the theorems they support}
\label{app:checks}

Scripts (directory \texttt{symbolic/} of the repository), each printing
\texttt{ALL PASS}:

\begin{center}
\small
\begin{tabular}{l l}
\hline
Check & Supports\\
\hline
ID1--ID5 (\texttt{verify\_interaction\_deformation.py}) &
Theorem~\ref{thm:first}\\
ID6--ID8 & Theorem~\ref{thm:secondgen}\\
ID9a--b & Theorem~\ref{thm:threeone}\\
ID10 & Theorem~\ref{thm:jordanscaling} ($R_2$)\\
SR1--SR3 (\texttt{verify\_interaction\_order3.py}) &
Theorem~\ref{thm:selection}\\
O3a--O3d & Theorem~\ref{thm:third}\\
O3e & Theorem~\ref{thm:jordanscaling} ($R_3$)\\
PT1--PT3 (\texttt{verify\_pt\_breaking.py}) &
Theorem~\ref{thm:complexpair}\\
PT4 & broken-PT tongues\\
PS1--PS2 & Proposition~\ref{prop:kallen}\\
PS-A--PS-G (\texttt{verify\_perfect\_square.py}) &
Theorem~\ref{thm:firstfield}\\
PS-H & Jordan-chain lemma used in
Proposition~\ref{prop:sectors}\\
TF1--TF7 (\texttt{verify\_two\_field.py}) &
Theorem~\ref{thm:twofield}, Proposition~\ref{prop:sectors}\\
SO1--SO7 (\texttt{verify\_sector\_obstruction.py}) &
Theorem~\ref{thm:sector} (numerics)\\
HX1 (\texttt{verify\_hardening.py}) &
Theorem~\ref{thm:sector} (exact value)\\
HX2--HX3 & Theorem~\ref{thm:confluent}\\
\hline
\end{tabular}
\end{center}

\end{document}
